% Bagian: sets-functions-relations
% Bab: size-of-sets
% Subbagian: equinumerous-sets

\documentclass[../../../include/open-logic-section]{subfiles}

\begin{document}

\olfileid[id]{sfr}{siz}{equ}

\olsection{Kesamaan Kardinalitas}

Kita mempunyai gagasan intuitif mengenai ``ukuran'' himpunan, yang bekerja
dengan baik untuk himpunan hingga. Namun, bagaimana dengan himpunan tak
hingga? Jika kita ingin merumuskan cara formal untuk membandingkan ukuran dua
himpunan dengan \emph{sebarang} ukuran, sebaiknya kita mulai dengan
mendefinisikan kapan dua himpunan berukuran sama. Berikut penjelasan Frege:
\begin{quote}
  Jika seorang pramusaji ingin memastikan bahwa ia telah meletakkan pisau
  tepat sebanyak piring di atas meja, ia tidak perlu menghitung salah satunya,
  asalkan ia meletakkan sebuah pisau di sebelah kanan setiap piring sehingga
  setiap pisau di atas meja berada di sebelah kanan suatu piring. Dengan
  demikian, piring dan pisau saling dikorelasikan secara tunggal, dan memang
  melalui hubungan spasial yang sama itu. \citep[\S70]{Frege1884}
\end{quote}
Gagasan dalam kutipan ini dapat dinyatakan melalui suatu definisi formal:

\begin{defn}\ollabel{comparisondef}
  $A$ \emph{berkardinalitas sama} dengan $B$, ditulis $\cardeq{A}{B}$, jika dan
  hanya jika terdapat !!a{bijection} $f \colon A \to B$.
\end{defn}

\begin{prop}\ollabel{equinumerosityisequi}
Kesamaan kardinalitas adalah relasi ekuivalensi.
\end{prop}

\begin{proof}
Kita harus menunjukkan bahwa kesamaan kardinalitas bersifat refleksif,
simetris, dan transitif. Misalkan $A, B$, dan $C$ adalah himpunan.

\emph{Refleksivitas.} Pemetaan identitas $\Id{A} \colon A \to A$, dengan
$\Id{A} (x) = x$ untuk semua $x \in A$, adalah !!a{bijection}. Jadi,
$\cardeq{A}{A}$.

\emph{Simetri.} Andaikan $\cardeq{A}{B}$, yakni terdapat !!a{bijection}
$f\colon A \to B$. Karena $f$ !!{bijective}, inversnya $f^{-1}$ ada dan juga
!!{bijective}. Jadi, $f^{-1}\colon B \to A$ adalah !!a{bijection}, sehingga
$\cardeq{B}{A}$.

\emph{Transitivitas.} Andaikan $\cardeq{A}{B}$ dan $\cardeq{B}{C}$, yakni
terdapat !!{bijection}s $f\colon A \to B$ dan $g\colon B \to C$. Komposisi
$\comp{f}{g}\colon A \to C$ adalah !!{bijective}, sehingga $\cardeq{A}{C}$.
\end{proof}

\begin{prop}
Jika $\cardeq{A}{B}$, maka $A$ adalah !!{enumerable} jika dan hanya jika $B$
demikian.
\end{prop}

\begin{editorial}
Bukti berikut menggunakan \olref[enm]{defn:enumerable} jika
\olref[enm]{sec} disertakan, dan \olref[enm-alt]{defn:enumerable} jika tidak.
\end{editorial}

\begin{proof}
Andaikan $\cardeq{A}{B}$, sehingga terdapat suatu !!{bijection} $f \colon A
\to B$, dan andaikan $A$ adalah !!{enumerable}.
\oliflabeldef{sfr:siz:enm:defn:enumerable}{
  Maka $A = \emptyset$ atau terdapat fungsi !!a{surjective} $g\colon \PosInt
  \to A$. Jika $A = \emptyset$, maka $B = \emptyset$ juga (jika tidak, akan
  terdapat !!a{element}~$y \in B$, tetapi tidak ada $x \in A$ dengan $f(x) =
  y$). Sebaliknya, jika $g\colon \PosInt \to A$ !!{surjective}, maka
  $\comp{g}{f} \colon \PosInt \to B$ !!{surjective}. Untuk melihatnya, ambil
  $y \in B$. Karena $f$ !!{surjective}, terdapat $x \in A$ sedemikian sehingga
  $f(x) = y$. Karena $g$ !!{surjective}, terdapat $n \in \PosInt$ sedemikian
  sehingga $g(n) = x$. Jadi,
\[
(\comp{g}{f})(n) = f(g(n)) = f(x) = y
\]
dan dengan demikian $\comp{g}{f}$ !!{surjective}. Jadi, $\comp{g}{f}$ adalah
enumerasi dari~$B$, sehingga $B$~adalah !!{enumerable}.}
{
  Maka $A = \emptyset$ atau terdapat !!a{bijection}~$g$ yang daerah hasilnya
  $A$ dan domainnya berupa $\Nat$ atau segmen awal bilangan asli. Jika $A =
  \emptyset$, maka $B = \emptyset$ juga (jika tidak, akan terdapat suatu~$y
  \in B$, tetapi tidak ada $x \in A$ sedemikian sehingga $f(x) = y$). Jadi,
  andaikan kita mempunyai !!{bijection}~$g$ tersebut. Maka $\comp{g}{f}$
  adalah !!a{bijection} dengan daerah hasil~$B$ dan domain yang sama dengan
  domain~$g$ (yakni $\Nat$ atau suatu segmen awalnya), sehingga $B$ adalah
  !!{enumerable}.}

Jika $B$ adalah !!{enumerable}, kita memperoleh bahwa $A$ juga
!!{enumerable} dengan mengulangi argumen menggunakan !!{bijection}
$f^{-1}\colon B \to A$ alih-alih~$f$.
\end{proof}

\begin{prob}
Tunjukkan bahwa jika $\cardeq{A}{C}$ dan $\cardeq{B}{D}$, serta $A \cap B = C
\cap D = \emptyset$, maka $\cardeq{A \cup B}{C \cup D}$.
\end{prob}

\begin{prob}
Tunjukkan bahwa jika $A$ tak hingga dan !!{enumerable}, maka
$\cardeq{A}{\Nat}$.
\end{prob}

\end{document}
